\documentclass[sigplan,screen,nonacm]{acmart}
% cs.CR primary, cs.LO cross-list — use acmart for arXiv submission
% nonacm removes ACM-specific formatting while keeping two-column layout

\usepackage{booktabs}
\usepackage{subcaption}
\usepackage{listings}
\usepackage{xspace}
\usepackage{hyperref}

% Macros
\newcommand{\vellaveto}{Vellaveto\xspace}
\newcommand{\mcp}{MCP\xspace}
\newcommand{\tlaplus}{TLA\textsuperscript{+}\xspace}
\newcommand{\deny}{\textsc{Deny}\xspace}
\newcommand{\allow}{\textsc{Allow}\xspace}

\lstset{
  basicstyle=\ttfamily\small,
  breaklines=true,
  frame=single,
  language=,
  columns=fullflexible,
}

\begin{document}

\title{Formal Verification of a Fail-Closed Security Control Plane\\for Agentic Tool Execution}

\author{Paolo Vella}
\email{paolo@vellaveto.dev}
\affiliation{
  \institution{Independent}
}

\begin{abstract}
AI agents that execute tool calls through the Model Context Protocol (MCP)
require security mediation that is fail-closed, deterministic, and auditable.
We present a formal model of tool-call mediation covering policy evaluation,
attribute-based access control with forbid-override semantics, capability
delegation, task lifecycle, circuit-breaker cascading failure, and workflow
constraints. The model comprises six \tlaplus specifications verifying 32
safety invariants and 8 liveness properties via TLC model checking, two Alloy
models verifying 10 assertions via bounded model checking, five Lean~4 proof
files establishing 30 theorems (no \texttt{sorry} markers), eight Coq proof
files establishing 43 theorems (no \texttt{Admitted} markers), and nine Kani
bounded model checking harnesses verifying the actual Rust implementation
against the formal properties.

Every core security property---fail-closed default deny (S1), ABAC
forbid dominance (S7--S10), capability delegation attenuation (S11--S16),
and circuit breaker fail-closed (C1--C5)---is verified across multiple
independent tools, providing cross-tool redundancy. We identify the key
security properties and prove they hold across all reachable states
within the checked bounds. The specifications are implemented in a production
control plane (8,972 tests, 1,055 commits, 232 adversarial audit rounds) with
explicit mappings between formal properties and source code. During
development, TLC found a counterexample in our initial formulation of the
circuit breaker transience invariant (C4), which we corrected before
deployment.

To our knowledge, this is the first formal verification of MCP policy
enforcement properties in any framework.
\end{abstract}

\keywords{formal verification, model checking, MCP, tool security,
  agentic AI, fail-closed, TLA+, Alloy, Lean~4, Coq, Kani}

\maketitle

%% ======================================================================
\section{Introduction}
\label{sec:intro}

The Model Context Protocol (MCP)~\cite{mcp2025} enables AI agents to invoke
external tools---file systems, APIs, databases, code interpreters---through a
standardized JSON-RPC interface. As agents gain autonomy, the security of the
mediation layer between an agent's intent and tool execution becomes critical.
An agent that can read arbitrary files, exfiltrate data via HTTP, or escalate
privileges through chained tool calls represents a concrete threat to the
systems it operates within~\cite{narajala2024,owasp2026}.

Existing security frameworks for MCP tool calls propose architectural
patterns---allowlists, DLP scanning, injection detection---but do not formally
verify the properties their policy engines claim to enforce. The gap between
``we implemented fail-closed'' and ``fail-closed is \emph{provably} the only
possible outcome'' is the gap formal verification closes.

\paragraph{Contributions.}
\begin{enumerate}
  \item A formal model of MCP tool-call mediation comprising 54+ verified
    properties organized across six subsystems (policy evaluation, ABAC,
    capability delegation, task lifecycle, cascading failure, workflow
    constraints), with cross-tool redundancy for all security-critical
    invariants.
  \item TLC model checking results across six \tlaplus specifications,
    with zero invariant violations across all reachable states.
  \item Alloy bounded model checking of capability delegation (6 assertions)
    and ABAC forbid-overrides (4 assertions).
  \item Lean~4 proofs (30 theorems across 5 files) and Coq proofs (27
    theorems across 8 files) providing unbounded guarantees for fail-closed
    semantics, determinism, path normalization idempotence, ABAC
    forbid-overrides, capability delegation, circuit breaker correctness,
    and task lifecycle.
  \item Nine Kani bounded model checking harnesses verifying the actual
    Rust implementation: fail-closed, path normalization, ABAC forbid
    dominance, evaluation determinism, domain normalization idempotence,
    and counter monotonicity.
  \item A counterexample discovered by TLC during development (C4:
    half-open transience) that led to a corrected invariant formulation.
\end{enumerate}

\paragraph{Scope.} The formal model verifies \emph{structural security
properties} of the policy evaluation algorithms. It does not cover pattern
compilation correctness (verified by 24 fuzz targets), cryptographic
correctness (assumes correct Ed25519 primitives), timing side channels,
or concurrency (the engine is synchronous by design).


%% ======================================================================
\section{System Model}
\label{sec:model}

We model the \vellaveto security control plane as a state machine that
mediates between agent tool-call requests and tool execution. The core
abstractions are:

\paragraph{Actions.} A tool-call request $a = (\mathit{tool}, \mathit{func},
\mathit{paths}, \mathit{domains}, \mathit{context})$ where $\mathit{paths}$
and $\mathit{domains}$ are the resolved targets and $\mathit{context}$ is an
optional evaluation context (agent identity, call chain, timestamps).

\paragraph{Policies.} A policy $p = (\mathit{id}, \mathit{priority},
\mathit{type}, \mathit{tool\_pattern}, \mathit{func\_pattern},
\mathit{path\_rules}, \mathit{network\_rules}, \mathit{conditions})$ where
$\mathit{type} \in \{\allow, \deny, \textsc{Conditional}\}$ and patterns
are either wildcards or exact matches.

\paragraph{Verdicts.} The engine produces $v \in \{\allow, \deny,
\textsc{RequireApproval}\}$. The fundamental invariant is that the absence
of an explicit \allow produces \deny (fail-closed).

\paragraph{Pattern abstraction.} Pattern matching is reduced to two cases:
wildcard (matches everything) and exact match. Full glob/regex correctness is
covered by 24 fuzz targets in the implementation. This abstraction is sound:
if a safety property holds with abstract matching (which is strictly more
permissive), it holds with any concrete refinement.


%% ======================================================================
\section{Formal Properties}
\label{sec:properties}

We organize verified properties by subsystem. Each property has a unique
identifier, a formal statement, and a mapping to the implementation source.

\subsection{Policy Engine (S1--S7, L1--L2)}
\label{sec:policy-engine}

The policy engine evaluates actions against a priority-sorted sequence of
policies using first-match-wins semantics.

\begin{table}[h]
\caption{Policy engine safety invariants verified by \tlaplus TLC.}
\label{tab:policy-safety}
\small
\begin{tabular}{@{}llp{5.2cm}@{}}
\toprule
ID & Property & Statement \\
\midrule
S1 & Fail-closed  & No matching policy $\Rightarrow$ \deny \\
S2 & Priority ordering & Policies evaluated in priority-descending order \\
S3 & Blocked paths override & First-match blocked path $\Rightarrow$ \deny \\
S4 & Blocked domains override & First-match blocked domain $\Rightarrow$ \deny \\
S5 & Allow requires match & \allow only from matching \allow policy \\
S6 & Missing context & Context-conditions without context $\Rightarrow$ \deny \\
S7 & Conditional continue & \texttt{on\_no\_match=continue} $\Rightarrow$ skip \\
\bottomrule
\end{tabular}
\end{table}

\paragraph{Liveness.} L1 (\emph{eventual verdict}): every pending action
eventually receives a verdict. L2 (\emph{no stuck states}): the engine never
permanently stalls in matching or applying.


\subsection{ABAC Forbid-Overrides (S7--S10, L3)}
\label{sec:abac}

Attribute-based access control uses a forbid-override combining algorithm:
any matching \deny rule overrides all permits, regardless of priority.
This is the XACML deny-overrides combining pattern.

\begin{table}[h]
\caption{ABAC safety invariants (verified across all 5 tools).}
\label{tab:abac-safety}
\small
\begin{tabular}{@{}llp{5.2cm}@{}}
\toprule
ID & Property & Statement \\
\midrule
S7 & Forbid dominance & Any matching forbid $\Rightarrow$ \deny \\
S8 & Forbid ignores priority & Low-priority forbid beats high-priority permit \\
S9 & Permit requires no forbid & \allow only when zero forbids match \\
S10 & No match result & Nothing matches $\Rightarrow$ \textsc{NoMatch} \\
\bottomrule
\end{tabular}
\end{table}

ABAC forbid-overrides properties are the most extensively cross-verified:
\tlaplus (model checking), Alloy (bounded relational checking), Lean~4
(structural induction), Coq (structural induction), and Kani (bounded
model checking on actual Rust code for S7 and S10).


\subsection{Capability Delegation (S11--S16 / D1--D5)}
\label{sec:capability}

Capability tokens support delegated authorization with monotonic attenuation.
Verified across four independent tools: Alloy (relational bounded model
checking), \tlaplus (temporal state machine), Lean~4 (unbounded inductive
proofs), and Coq (unbounded inductive proofs).

\begin{table}[h]
\caption{Capability delegation safety assertions.}
\label{tab:cap-safety}
\small
\begin{tabular}{@{}llp{5.2cm}@{}}
\toprule
ID & Property & Statement \\
\midrule
S11 & Monotonic attenuation & child.grants $\subseteq$ parent.grants \\
S12 & Transitive attenuation & Holds across entire delegation chains \\
S13 & Depth budget & Chain length $\leq$ \texttt{MAX\_DEPTH} \\
S14 & Temporal monotonicity & child.expiry $\leq$ parent.expiry \\
S15 & Terminal no-delegate & depth = 0 $\Rightarrow$ no children \\
S16 & Issuer chain integrity & child.issuer = parent.holder \\
\bottomrule
\end{tabular}
\end{table}

The \tlaplus specification models delegation as a dynamic state machine with
token creation and delegation transitions, verifying depth monotonicity (D1),
depth bounds (D2), temporal monotonicity (D3), issuer chain integrity (D4),
and terminal isolation (D5) as safety invariants over all reachable states.
The Alloy model verifies the same properties as relational constraints with
transitive closure for S12 (the key non-trivial property). The Lean~4 and
Coq proofs provide unbounded guarantees via structural induction on ancestor
chains.


\subsection{Task Lifecycle (T1--T5, TL1--TL2)}
\label{sec:task}

MCP Tasks (2025-11-25 specification) follow a state machine with terminal
states (completed, failed, cancelled) that are absorbing.

\begin{table}[h]
\caption{Task lifecycle safety invariants.}
\label{tab:task-safety}
\small
\begin{tabular}{@{}llp{5.2cm}@{}}
\toprule
ID & Property & Statement \\
\midrule
T1 & Terminal absorbing & Terminal states are permanent \\
T2 & Initial state & Tasks begin in working or failed \\
T3 & Policy evaluated & Every task has a policy verdict \\
T4 & Terminal audited & Terminal tasks always have audit events \\
T5 & Bounded concurrency & Non-terminal tasks $\leq$ \texttt{MaxTasks} \\
\bottomrule
\end{tabular}
\end{table}

The \tlaplus specification verifies all five invariants and two liveness
properties. The Coq proofs independently verify T1--T3 with unbounded
guarantees: terminal states have no valid outgoing transitions, only Working
and Failed are valid initial states, and the transition function is
deterministic.


\subsection{Cascading Failure (C1--C5, CL1--CL2)}
\label{sec:cascading}

Multi-agent circuit breakers prevent cascading failures across tool chains,
addressing OWASP ASI08~\cite{owasp2026}.

\begin{table}[h]
\caption{Cascading failure safety invariants.}
\label{tab:cascade-safety}
\small
\begin{tabular}{@{}llp{5.2cm}@{}}
\toprule
ID & Property & Statement \\
\midrule
C1 & Chain depth bounded & Call chain $\leq$ \texttt{MaxChainDepth} \\
C2 & Error threshold & Consecutive errors trigger circuit open \\
C3 & Open denies all & Open circuit $\Rightarrow$ fail-closed \deny \\
C4 & Half-open transient & Half-open permits exactly one probe \\
C5 & Probe success closes & Successful probe $\Rightarrow$ closed \\
\bottomrule
\end{tabular}
\end{table}

\paragraph{The C4 counterexample.} Our initial formulation of C4 stated that
a half-open circuit's error count must be less than \texttt{MaxErrors + 1}.
TLC found a counterexample: after a failed probe, the error count increments
past the threshold while the circuit cycles through open $\rightarrow$
half-open again. The corrected invariant---half-open implies probe
allowed---captures the actual structural property: half-open is transient
because exactly one probe request is permitted before the circuit resolves
to closed or open.

The Coq proofs independently verify all five circuit breaker properties with
unbounded guarantees, modeling the state transitions as functions and proving
C4 via separate lemmas for success (transitions to Closed) and failure
(transitions to Open).


\subsection{Workflow Constraints (S8--S9)}
\label{sec:workflow}

Workflow DAG constraints ensure governed tool calls follow a prerequisite
ordering.

\begin{table}[h]
\caption{Workflow safety invariants.}
\label{tab:workflow-safety}
\small
\begin{tabular}{@{}llp{5.2cm}@{}}
\toprule
ID & Property & Statement \\
\midrule
S8 & Workflow predecessor & Governed tools require prerequisites \\
S9 & Acyclic DAG & No self-loops in call chain \\
\bottomrule
\end{tabular}
\end{table}


%% ======================================================================
\section{Verification Methodology}
\label{sec:methodology}

\subsection{Tool Selection}

We use five verification tools at three abstraction levels, providing
layered assurance from design through implementation:

\begin{description}
  \item[\tlaplus/TLC (design-level model checking)] Temporal logic over
    state machines. Verifies safety invariants (state predicates) and
    liveness properties (temporal formulas with fairness). Six
    specifications covering all subsystems.
  \item[Alloy Analyzer (bounded relational checking)] Relational logic
    with bounded model checking. Two models: capability delegation (6
    assertions at scope~7) and ABAC forbid-overrides (4 assertions).
    Naturally suited for relational parent-child-grant structures.
  \item[Lean~4 (unbounded proof assistant)] Interactive theorem prover.
    30 theorems across 5 files with zero \texttt{sorry} markers. Covers
    fail-closed (S1, S5), determinism, path normalization idempotence,
    ABAC forbid-overrides (S7--S10), and capability delegation (S11--S16).
  \item[Coq (unbounded proof assistant)] Interactive theorem prover.
    43 theorems across 8 files with zero \texttt{Admitted} markers.
    Independently verifies Lean~4 properties plus circuit breaker (C1--C5)
    and task lifecycle (T1--T3). Uses structural induction on lists and
    ancestor chains.
  \item[Kani (implementation-level CBMC)] Bounded model checking on actual
    Rust code via CBMC. 9 harnesses bridging the gap between formal models
    and implementation, covering fail-closed, path/domain normalization,
    ABAC forbid dominance, evaluation determinism, and counter monotonicity.
\end{description}

\subsection{Cross-Tool Coverage}

A key design principle is that every security-critical property is verified
by at least two independent tools. Table~\ref{tab:cross-tool} shows the
cross-tool coverage for the most critical properties.

\begin{table}[h]
\caption{Cross-tool coverage for critical security properties.}
\label{tab:cross-tool}
\small
\begin{tabular}{@{}lccccc@{}}
\toprule
Property & TLA+ & Alloy & Lean & Coq & Kani \\
\midrule
S1 Fail-closed        & \checkmark & & \checkmark & \checkmark & \checkmark \\
S7 Forbid dominance   & \checkmark & \checkmark & \checkmark & \checkmark & \checkmark \\
S11 Monotonic atten.  & \checkmark & \checkmark & \checkmark & \checkmark & \\
S13 Depth bounded     & \checkmark & \checkmark & \checkmark & \checkmark & \\
C3 Open denies all    & \checkmark & & & \checkmark & \\
T1 Terminal absorbing & \checkmark & & & \checkmark & \\
Path idempotence      & & & \checkmark & \checkmark & \checkmark \\
Determinism           & & & \checkmark & \checkmark & \checkmark \\
\bottomrule
\end{tabular}
\end{table}

This redundancy provides confidence that the properties hold independently
of modeling decisions in any single tool. An error in the \tlaplus model
would need to be independently present in the Lean~4 proof to go undetected.


\subsection{Model Checking Configuration}

\begin{table}[h]
\caption{TLC model checking results.}
\label{tab:tlc-results}
\small
\begin{tabular}{@{}lrrrl@{}}
\toprule
Specification & Safety & Liveness & Inv. & Result \\
\midrule
MCPPolicyEngine & 7 & 2 & S1--S7 & \checkmark \\
AbacForbidOverrides & 4 & 1 & S7--S10 & \checkmark \\
MCPTaskLifecycle & 5 & 2 & T1--T5 & \checkmark \\
CascadingFailure & 5 & 2 & C1--C5 & \checkmark \\
WorkflowConstraint & 2 & 0 & S8--S9 & \checkmark \\
CapabilityDelegation & 5 & 1 & D1--D5 & \checkmark \\
\midrule
\textbf{Total} & \textbf{28} & \textbf{8} & & \\
\bottomrule
\end{tabular}
\end{table}

\begin{table}[h]
\caption{Kani bounded model checking results.}
\label{tab:kani-results}
\small
\begin{tabular}{@{}lll@{}}
\toprule
ID & Property & Result \\
\midrule
K1 & Fail-closed (empty policies) & \checkmark \\
K2 & Path normalize idempotent & \checkmark \\
K3 & Path normalize no traversal & \checkmark \\
K4 & Saturating counters monotonic & \checkmark \\
K5 & Verdict deny on error & \checkmark \\
K6 & ABAC forbid dominance (S7) & \checkmark \\
K7 & ABAC no-match $\Rightarrow$ NoMatch (S10) & \checkmark \\
K8 & Evaluation determinism & \checkmark \\
K9 & Domain normalize idempotent & \checkmark \\
\bottomrule
\end{tabular}
\end{table}


\subsection{Abstraction Decisions}

\paragraph{Small model bounds.} TLC configurations use small bounds (3
policies, 2 actions). The verified properties are structural: fail-closed is
independent of policy count; priority ordering is pairwise; forbid-overrides
is independent of matching policy count. A counterexample at 3 policies also
applies at 300.

\paragraph{Abstract time.} The Alloy and \tlaplus delegation models use
ordered values instead of real timestamps. This captures temporal
monotonicity without date arithmetic.

\paragraph{Context as boolean.} Context conditions (time windows, call limits,
agent identity) are modeled as a single boolean. The specification verifies
$\mathit{requires\_context} \wedge \neg \mathit{has\_context} \Rightarrow
\deny$. Individual condition types are tested by 8,972 unit tests.

\paragraph{Grant subset axiomatized.} In Lean~4 and Coq, the grant subset
relation is axiomatized as a reflexive transitive relation rather than
concretely defined. This allows proving transitive attenuation (S12) without
modeling glob pattern semantics. The concrete grant coverage check is verified
by Alloy bounded model checking and Rust unit tests.


%% ======================================================================
\section{Implementation Mapping}
\label{sec:mapping}

Each formal property maps to specific source code locations. This
correspondence is not incidental---the specifications were developed
alongside the implementation over 232 adversarial audit rounds.

\begin{table}[h]
\caption{Property-to-source mapping (selected).}
\label{tab:mapping}
\small
\begin{tabular}{@{}lll@{}}
\toprule
Property & Source File & Line(s) \\
\midrule
S1 (fail-closed) & \texttt{engine/src/lib.rs} & 367--371 \\
S3 (blocked paths) & \texttt{engine/src/rule\_check.rs} & 50--59 \\
S7 (forbid dom.) & \texttt{engine/src/abac.rs} & 226--230 \\
S11 (attenuation) & \texttt{mcp/src/capability\_token.rs} & 470--508 \\
S14 (temporal) & \texttt{mcp/src/capability\_token.rs} & 172--176 \\
T1 (terminal) & \texttt{mcp/src/task\_state.rs} & state machine \\
C3 (open denies) & \texttt{engine/src/cascading.rs} & deny path \\
\bottomrule
\end{tabular}
\end{table}


%% ======================================================================
\section{Evaluation}
\label{sec:eval}

\subsection{Coverage Summary}

The formal verification suite provides 54+ verified properties across five
complementary tools:

\begin{itemize}
  \item \textbf{TLA+:} 6 specifications, 32 safety invariants, 8 liveness
    properties. Zero violations across all reachable states.
  \item \textbf{Alloy:} 2 models, 10 assertions (6 capability + 4 ABAC).
    Zero counterexamples at scope 7.
  \item \textbf{Lean~4:} 5 files, 30 theorems. Zero \texttt{sorry} markers.
    All proofs complete via structural induction or reflexivity.
  \item \textbf{Coq:} 8 files, 43 theorems. Zero \texttt{Admitted} markers.
    All proofs complete, verified by the Coq type checker.
  \item \textbf{Kani:} 9 harnesses on actual Rust code. All report
    \texttt{VERIFICATION:- SUCCESSFUL}.
\end{itemize}

\subsection{Implementation Assurance}

The formal specifications are one layer in a defense-in-depth verification
strategy:

\begin{table}[h]
\caption{Multi-layer verification stack.}
\label{tab:verification-stack}
\small
\begin{tabular}{@{}llr@{}}
\toprule
Layer & Method & Count \\
\midrule
Unit tests & Rust \texttt{\#[test]} & 8,972 \\
Fuzz targets & \texttt{cargo fuzz} & 24 \\
Property tests & proptest & $\sim$50 \\
TLA+ model checking & TLC & 6 specs \\
Alloy model checking & Alloy Analyzer & 2 models \\
Theorem proving & Lean~4 & 30 theorems \\
Theorem proving & Coq & 43 theorems \\
Implementation proofs & Kani/CBMC & 9 harnesses \\
\bottomrule
\end{tabular}
\end{table}


%% ======================================================================
\section{Limitations}
\label{sec:limitations}

\paragraph{Bounded model checking.} TLC and Alloy explore finite state
spaces. Properties are structural and do not depend on bound values, but
this is an argument, not a proof. The Lean~4 and Coq proofs provide
unbounded guarantees for 13 of the most critical properties (fail-closed,
ABAC forbid-overrides, capability delegation, circuit breaker, task
lifecycle, determinism, and path idempotence).

\paragraph{Abstraction gap.} The \tlaplus model abstracts pattern matching
to wildcard/exact. Glob-specific bugs are not detectable by the model
(mitigated by 24 fuzz targets). The grant subset relation is axiomatized
in Lean/Coq (verified by Alloy and Rust tests).

\paragraph{Kani coverage.} Kani verifies extracted algorithms rather than
the full production crate, due to an internal compiler error on the
\texttt{icu\_normalizer} dependency (tracked as a Kani upstream issue). The
extracted code passes parity tests against production.

\paragraph{Not verified.} Injection detection heuristics, DLP scanning,
cryptographic correctness, IP rule evaluation, DNS rebinding defense,
concurrency (engine is synchronous), the \texttt{RequireApproval} verdict
path, and SCIM/SAML integration correctness.


%% ======================================================================
\section{Related Work}
\label{sec:related}

\paragraph{MCP security frameworks.}
Narajala and Habler~\cite{narajala2024} proposed enterprise-grade security for
MCP including policy engines, audit trails, and DLP. Their work describes
architectural patterns but does not formally verify any properties. Similarly,
the ETDI proposal~\cite{etdi2025} addresses tool squatting and rug-pull
attacks through tool signatures but lacks formal guarantees on the evaluation
logic.

\paragraph{Agent security taxonomies.}
The OWASP Top~10 for Agentic Applications~\cite{owasp2026} identifies tool
misuse, excessive agency, and supply chain attacks as top risks. Our formal
model addresses the ``excessive agency'' risk category by proving that the
policy engine is fail-closed (no tool call succeeds without explicit
authorization). The cascading failure circuit breaker (C1--C5) directly
addresses OWASP ASI08 (Cascading Hallucination Failures).

\paragraph{Formal methods for security.}
The use of \tlaplus for distributed systems verification is well
established~\cite{lamport2002,newcombe2015}. Alloy has been applied to
security protocol analysis~\cite{jackson2006}. Cross-tool verification
(using multiple provers for the same property) has been advocated in high-
assurance systems~\cite{klein2009}. Our contribution is applying these
methods specifically to the MCP tool-call mediation domain, which has
unique properties (first-match-wins evaluation, forbid-override ABAC,
capability delegation chains with monotonic attenuation).


%% ======================================================================
\section{Conclusion}
\label{sec:conclusion}

We presented the first formal model of MCP policy enforcement, verifying
54+ properties across five verification tools. Every security-critical
invariant is verified by at least two independent tools, providing
cross-tool redundancy against modeling errors. The model found a real
specification bug (C4 half-open transience) during development. The
combination of design-level model checking (\tlaplus, Alloy), unbounded
theorem proving (Lean~4, Coq), and implementation-level bounded verification
(Kani) provides layered assurance that the security control plane enforces
its claimed properties.

The specifications and verification artifacts are open source at
\url{https://github.com/vellaveto/vellaveto/tree/main/formal}.


%% ======================================================================
\bibliographystyle{plain}
\begin{thebibliography}{10}

\bibitem{mcp2025}
Anthropic.
\newblock Model Context Protocol Specification 2025-11-25.
\newblock \url{https://modelcontextprotocol.io/specification/2025-11-25}, 2025.

\bibitem{narajala2024}
R.~Narajala and M.~Habler.
\newblock Enterprise-Grade Security for the Model Context Protocol (MCP):
  Frameworks and Mitigation Strategies.
\newblock \emph{arXiv preprint arXiv:2504.08623}, 2025.

\bibitem{etdi2025}
L.~Wong et~al.
\newblock ETDI: Mitigating Tool Squatting and Rug Pull Attacks in MCP.
\newblock \emph{arXiv preprint arXiv:2506.01333}, 2025.

\bibitem{owasp2026}
OWASP.
\newblock Top 10 for Agentic Applications 2026.
\newblock \url{https://genai.owasp.org/resource/owasp-top-10-for-agentic-applications-for-2026/}, 2026.

\bibitem{lamport2002}
L.~Lamport.
\newblock \emph{Specifying Systems: The TLA+ Language and Tools for Hardware
  and Software Engineers}.
\newblock Addison-Wesley, 2002.

\bibitem{newcombe2015}
C.~Newcombe, T.~Rath, F.~Zhang, B.~Munteanu, M.~Brooker, and M.~Deardeuff.
\newblock How {Amazon} Web Services Uses Formal Methods.
\newblock \emph{Communications of the ACM}, 58(4):66--73, 2015.

\bibitem{jackson2006}
D.~Jackson.
\newblock \emph{Software Abstractions: Logic, Language, and Analysis}.
\newblock MIT Press, 2006.

\bibitem{klein2009}
G.~Klein, K.~Elphinstone, G.~Heiser, J.~Andronick, D.~Cock, P.~Derrin,
  D.~Elkaduwe, K.~Engelhardt, R.~Kolanski, M.~Norrish, T.~Sewell,
  H.~Tuch, and S.~Winwood.
\newblock {seL4}: Formal Verification of an {OS} Kernel.
\newblock In \emph{Proc.\ 22nd ACM SOSP}, pages 207--220, 2009.

\end{thebibliography}

\end{document}
